% ============================================================
% Supplementary File S1: Search Strategy and Review Protocol
% ============================================================
% NOTE TO AUTHOR (updated June 4, 2026): The a priori protocol below is
% complete — criteria, concept groups, and full field-tagged query strings
% are the real, defined protocol. WHAT REMAINS: the formal database searches
% have not yet been EXECUTED. Once you run them, record (a) the execution
% date per database, and (b) the per-database record counts, in the cells
% marked [[TO RECORD: ...]], then refresh the PRISMA counts in the
% manuscript and remove the DRAFT banner. Do NOT submit before execution.
% ============================================================

\newpage
\section*{Supplementary File S1: Search Strategy and Review Protocol}
\label{sec:S1}

\noindent\colorbox{lowred}{\parbox{0.97\linewidth}{\textcolor{lowredtext}{\textbf{%
DRAFT --- not for submission as-is.} The protocol below (criteria, concept
groups, and full query strings) is final and was defined a priori. The formal
database searches have \textbf{not yet been executed}: the current evidence base
is a provisional seed set. Before submission, execute the searches, record the
per-database counts and execution dates in the cells marked
\texttt{[[TO RECORD: \ldots]]}, refresh the PRISMA counts throughout the
manuscript, and remove this banner.}}}

\bigskip

\noindent\textbf{Purpose:} Document the protocol that was defined prior to search
execution---databases, date coverage, full Boolean query strings, inclusion and
exclusion criteria, and the U-T-I-O coding framework---so that the systematic
review is reproducible.

\subsection*{1. Databases and Coverage}

\noindent\textbf{Primary databases:} PubMed/MEDLINE, IEEE Xplore, Scopus. \\
\textbf{Supplementary sources:} Web of Science, CINAHL, Google Scholar,
hand-searching of reference lists, and expert recommendations. \\
\textbf{Publication window:} 2003 through May 2026. The empirical search applies
a 2015 lower date limit (HITECH incentive era through current FHIR mandate
implementation); foundational theory published from 2003 onwards (e.g., DeLone
\& McLean) is retained as background rather than retrieved through the empirical
search. \\
\textbf{Search executed:} [[TO RECORD: date(s) each database was searched]].

\subsection*{2. Concept Groups (Search Terms)}

The search combined two concept groups with the Boolean operator \texttt{AND};
terms within each group were combined with \texttt{OR}.

\renewcommand{\arraystretch}{1.3}
\begin{table}[H]
\footnotesize
\centering
\begin{tabular}{p{3.0cm} p{11.0cm}}
\toprule
\rowcolor{headerblue}
\textcolor{white}{\textbf{Concept group}} &
\textcolor{white}{\textbf{Terms (OR-combined)}} \\
\midrule
\rowcolor{rowgray}
Technology domain &
``electronic health record,'' ``EHR,'' ``electronic medical record,'' ``EMR,''
``health information exchange,'' ``HIE,'' ``eHealth,'' ``health information
technology,'' ``interoperability'' \\
Adoption domain &
``adoption,'' ``implementation,'' ``barriers,'' ``facilitators,''
``determinants,'' ``success factors,'' ``uptake'' \\
\bottomrule
\end{tabular}
\end{table}

\subsection*{3. Full Boolean Query Strings (per database)}

\noindent\textit{The field-tagged query strings below were executed as recorded.
For PubMed/MEDLINE, both concept groups were restricted to the Title field to
yield a focused, single-reviewer--screenable corpus. This is a deliberate
precision-over-recall scoping decision appropriate to a single-author focused
review: it reliably captures studies whose primary subject is eHealth-systems
adoption (signalled in the title) at the cost of omitting studies that treat
adoption only peripherally in the abstract or body. This recall limitation is
acknowledged in the manuscript Limitations section.}

\medskip
\noindent\textbf{PubMed/MEDLINE} (field: Title; publication years 2015--2026;
article type unrestricted; English-language eligibility applied at screening): \\
{\footnotesize\begin{verbatim}
("electronic health record"[Title] OR "EHR"[Title]
OR "electronic medical record*"[Title]
OR "health information exchange"[Title] OR "HIE"[Title]
OR "eHealth"[Title] OR "e-health"[Title]
OR "digital health"[Title]
OR "health information technology"[Title]
OR "interoperability"[Title])
AND
("adoption"[Title] OR "implementation"[Title]
OR "acceptance"[Title] OR "barrier*"[Title]
OR "facilitator*"[Title] OR "determinant*"[Title]
OR "success factor*"[Title] OR "uptake"[Title])
AND ("2015/01/01"[PDAT] : "2026/12/31"[PDAT])
\end{verbatim}}

\medskip
\noindent\textbf{IEEE Xplore} (Command Search; content type: Conferences and
Journals; publication years 2015--2026): \\
{\footnotesize\begin{verbatim}
("Document Title":"electronic health record" OR "Document Title":"EHR"
OR "Document Title":"electronic medical record"
OR "Document Title":"health information exchange" OR "Document Title":"HIE"
OR "Document Title":"eHealth" OR "Document Title":"e-health"
OR "Document Title":"digital health"
OR "Document Title":"health information technology"
OR "Document Title":"interoperability")
AND
("Document Title":"adoption" OR "Document Title":"implementation"
OR "Document Title":"acceptance" OR "Document Title":"barrier"
OR "Document Title":"facilitator" OR "Document Title":"determinant"
OR "Document Title":"success factor" OR "Document Title":"uptake")
\end{verbatim}}

\medskip
\noindent\textbf{Scopus} (document types: Article, Review; English): \\
{\footnotesize\begin{verbatim}
TITLE-ABS-KEY ( "electronic health record" OR "EHR"
OR "electronic medical record" OR "health information exchange" OR "HIE"
OR "eHealth" OR "e-health" OR "digital health"
OR "health information technology" OR "interoperability" )
AND TITLE-ABS-KEY ( "adoption" OR "implementation" OR "acceptance"
OR "barriers" OR "facilitators" OR "determinants" OR "success factors"
OR "uptake" )
AND PUBYEAR > 2014 AND PUBYEAR < 2027
AND DOCTYPE ( ar OR re ) AND LANGUAGE ( english )
\end{verbatim}}

\medskip
\noindent\textbf{Web of Science} (Core Collection; 2015 onwards; English; Article or Review): \\
{\footnotesize\begin{verbatim}
TS=("electronic health record" OR "EHR" OR "health information exchange"
OR "eHealth" OR "digital health" OR "interoperability")
AND TS=("adoption" OR "implementation" OR "barrier*" OR "facilitator*"
OR "determinant*")
\end{verbatim}}

\medskip
\noindent\textbf{CINAHL} (2015 onwards; English; Peer Reviewed): \\
{\footnotesize\begin{verbatim}
(MH "Medical Informatics" OR TI eHealth OR TI "EHR"
OR TI "electronic health record")
AND (TI adoption OR TI implementation OR TI barrier OR TI facilitator
OR TI determinant)
\end{verbatim}}

\medskip
\noindent\textbf{Google Scholar} (supplementary; first 5 pages by relevance,
manually de-duplicated against primary databases): \\
{\footnotesize\begin{verbatim}
"eHealth adoption" OR "EHR implementation barriers"
OR "health information exchange determinants"
OR "digital health adoption framework"
\end{verbatim}}

\medskip
\noindent\textbf{Hand-search and expert recommendations} (supplementary):
hand-search of reference lists from retained full-text sources; citations
recommended by domain advisors.

\subsection*{4. Per-Database Record Counts}

\noindent\textit{The formal database searches were executed on June 4, 2026.
PubMed/MEDLINE and IEEE Xplore were searched directly; Google Scholar was run as
a supplementary recall check. Scopus, Web of Science, and CINAHL could not be
searched, as institutional subscription access was unavailable; their omission
is an acknowledged coverage limitation (see manuscript Limitations). The Google
Scholar sweep returned no unique within-window records: every relevant 2015--2026
result it surfaced was already retrieved by the PubMed/MEDLINE search, which
provides evidence of adequate PubMed recall. Three pre-2015 foundational reviews
surfaced by Google Scholar fall outside the empirical search window; one (Vest,
2010, on the TOE framework) is retained as a background/theoretical citation and
is not counted in the screened corpus.}

\renewcommand{\arraystretch}{1.25}
\begin{table}[H]
\footnotesize
\centering
\begin{tabular}{p{6.5cm} >{\centering\arraybackslash}p{3.0cm} >{\centering\arraybackslash}p{2.8cm}}
\toprule
\rowcolor{headerblue}
\textcolor{white}{\textbf{Source}} &
\textcolor{white}{\textbf{Records identified}} &
\textcolor{white}{\textbf{Date searched}} \\
\midrule
\rowcolor{rowgray}
PubMed/MEDLINE & 1{,}591 & 2026-06-04 \\
IEEE Xplore & 97 & 2026-06-04 \\
Google Scholar (supplementary) & 0 unique & 2026-06-04 \\
\rowcolor{rowgray}
Scopus & not searched (no access) & --- \\
Web of Science & not searched (no access) & --- \\
\rowcolor{rowgray}
CINAHL & not searched (no access) & --- \\
\midrule
\textbf{Total identified} & \textbf{1{,}688} & --- \\
\textbf{Duplicates removed} & \textbf{22} & --- \\
\textbf{After duplicates removed (screened)} & \textbf{1{,}666} & --- \\
\bottomrule
\end{tabular}
\end{table}

\noindent\textit{De-duplication was performed in Rayyan (rayyan.ai) using
automated exact-match detection (title, authors, journal, year, DOI; minimum
overall similarity 97\%). Of 66 candidate pairs flagged, 22 confirmed duplicates
were removed and the remainder retained, leaving 1{,}666 unique records for
title-and-abstract screening.}

\subsection*{5. Inclusion Criteria}

Records were included if they:
\begin{enumerate}
\item appeared in peer-reviewed journals or as institutional/government reports;
\item focused on eHealth, EHR, EMR, HIE, health information technology, or clinical
  interoperability;
\item addressed adoption, implementation, barriers, facilitators, or determinants of
  system uptake; and
\item employed empirical, systematic review, scoping review, or framework-based
  research designs.
\end{enumerate}

\subsection*{6. Exclusion Criteria}

Records were excluded if they:
\begin{enumerate}
\item were opinion editorials, letters, or conference abstracts without full
  methodology;
\item addressed non-healthcare IT systems without healthcare application;
\item were published in languages other than English; or
\item involved the primary collection of identifiable human-subject data.
\end{enumerate}

\subsection*{7. Screening and Data-Extraction Procedure}

Retrieved records were deduplicated using reference-management software. Title and
abstract screening was applied to all retained records using the criteria above.
Full-text review was conducted for records passing the initial screen. Data
extraction applied the U-T-I-O coding framework (construct definitions and example
codes in the manuscript Coding Schema table; case-level scoring rubric in
Supplementary File~S3). Reasons for full-text exclusion are itemised in
Supplementary File~S2.

\medskip
\noindent\textbf{Protocol registration:} Prospective protocol registration was not
undertaken: the review's outcomes are implementation and adoption determinants
rather than patient health outcomes, placing it outside PROSPERO's registration
scope, and the synthesis draws exclusively on published literature and public
documentation. The protocol documented here was defined a priori, prior to
search execution.

\medskip
\noindent\textit{Last updated: June 4, 2026 (protocol final; formal search
execution pending)}
